\chapter{Source file overview}
\label{app:source}

The prover sources live in \texttt{\textasciitilde/prover/}.  This appendix is a
map of the tree, not a design document: it says what each file \emph{is}, so a
reader following a cross-reference or extending the system can find the right
place.  The layout is:

\begin{itemize}
\item \textbf{Root} (\texttt{\textasciitilde/prover/*.scm}) --- the proof
  kernel, the reader and library record, the interactive command layer, and the
  documentation, proof-printing, and proof-driving subsystems.
\item \textbf{\texttt{structure-library/}} --- definitions, structures,
  vocabulary, and warranted supports: the mathematical objects the system knows
  at startup.
\item \textbf{\texttt{theorem-library/}} --- proofs that reach \texttt{(qed
  \ldots)}, loaded, gated, and counted.
\item \textbf{\texttt{calculus/}} --- probes and stress tests.  \emph{Mostly}
  outside the load path, but not wholly: \texttt{load.scm} names
  \texttt{calculus/compact-tb-proof},
  \texttt{calculus/compact-complete-proof} and
  \texttt{calculus/finite-ball-subcover-proof}, and it gives every
  \texttt{calculus/} file the same fresh-environment containment it gives
  \texttt{theorem-library/} precisely because such a file can be loaded.  An
  edit here is therefore not automatically inert.
\item \textbf{\texttt{reference/}} --- for the most part generated at every
  load (\texttt{PSS.md}, \texttt{THEOREMS.md}, \texttt{DEFINITIONS.md},
  \texttt{PROOF-DEBT.md}, \texttt{TACTICS.md}, \ldots), and those files must
  not be edited by hand; \texttt{LIBRARY.md}, \texttt{KERNEL.md},
  \texttt{KERNEL-RULES.md}, \texttt{REVIEW.md}, \texttt{USABILITY-REVIEW.md}
  and \texttt{VNB-TEST.md} are written by hand.
\item \textbf{\texttt{certificates/}} --- the proof-certificate store, one
  file per proof file.
\item \textbf{\texttt{libraries.sexp}} --- the assignment of files to
  libraries.
\item \textbf{\texttt{emacs/}} --- the workspace.
\item \textbf{\texttt{preambles/}} --- the preamble rule files.
\item \textbf{\texttt{notation/}} --- the printer tables (\texttt{base.tbl} and others).
\item \textbf{\texttt{docs/}} --- this manual and the design notes.
\item \textbf{\texttt{archive/}} --- retired and spliced originals.
\end{itemize}

\begin{remark}
The file listings below are \emph{generated} from the tree by
\texttt{docs/gen-source-appendix.py}, which harvests each file's
\texttt{;;;\ \textit{name} -- \textit{description}} header.  Regenerate after
adding or renaming files so this appendix cannot drift from the sources it
describes:
\begin{VNBsnip}
python3 docs/gen-source-appendix.py
\end{VNBsnip}
\end{remark}

%% AUTO-GENERATED by docs/gen-source-appendix.py -- do not edit by hand.
%% Regenerate:  python3 docs/gen-source-appendix.py
%% Descriptions are harvested from each file's `;;; name -- desc' header.

\section{Kernel and subsystems (\texttt{\textasciitilde/prover/})}
\label{sec:src-root}
The 82 Scheme files at the root of the tree: the proof kernel, the reader, the theory record, the interactive command layer, and the documentation / proof-printing / driver subsystems.

\begin{longtable}{@{}p{0.30\textwidth}p{0.63\textwidth}@{}}
\toprule
\textbf{File} & \textbf{Description} \\
\midrule
\endhead
\texttt{all-theorems.scm} & THE BOOK OF ALL PROVEN THEOREMS (the user's roadmap, \\
\texttt{arith-eval.scm} & ground arithmetic decision procedure \\
\texttt{audit.scm} & library hygiene diagnostics \\
\texttt{band-dump.scm} & dump what a saved band holds, one deterministic line per name \\
\texttt{calc.scm} & the ``directive chain checker'' (notes-27): ground a goal by a \\
\texttt{certificates.scm} & PROOF CERTIFICATES: a load that runs no proof \\
\texttt{cite-index.scm} & MEASURED citation affinity: P(theorem cited | head in goal) \\
\texttt{clobber-guard.scm} & the gate CLAUDE.md says does not exist \\
\texttt{contexts.scm} & make-wff, the interactive entry point for building a <wff> \\
\texttt{contra.scm} & (contra): close a branch whose CONTEXT is arithmetically \\
\texttt{core-capture.scm} & CORE REPLAY, steps 1-3: capture, the kernel bill, the lowering \\
\texttt{core-replay-checks.scm} & the suite checks of the CORE REPLAYER (core-replay.scm), loaded by \\
\texttt{core-replay.scm} & the CORE REPLAYER (core replay, design step 4) and the \\
\texttt{counterexample.scm} & probes that try to REFUTE the focus goal \\
\texttt{deduction-graphs.scm} & deduction graph: sequent nodes and inference nodes \\
\texttt{driver-kit.scm} & the proof-driving helpers that more than one file needs \\
\texttt{equal-patch.scm} & MIT/GNU Scheme 12.1's equal? without its per-call hash table \\
\texttt{errors.scm} & VNB error value and guard wrapper \\
\texttt{expressions.scm} & FOL expressions over VNB-style set theory \\
\texttt{extend-band.scm} & extend a saved band instead of rebuilding it (2026-09-23) \\
\texttt{functoid-gate.scm} & the def-functoid DEFINITION LOG and the load-end gate on \\
\texttt{glossary.scm} & THE ALPHABETICAL GLOSSARY: every name in VNB, in one list \\
\texttt{harvest.scm} & MEASURE the What Now panel against the library's own proofs \\
\texttt{ineq-supply.scm} & the What Now lane that SUPPLIES the inequalities the \\
\texttt{input-context.scm} & generalized context-discharge + operator-remap engine \\
\texttt{interactive.scm} & short-form commands for interactive proof sessions \\
\texttt{kernel-map-static.scm} & INSTRUMENT, not library \\
\texttt{kernel-map-trace.scm} & INSTRUMENT, not library \\
\texttt{libraries.scm} & THE LIBRARIES (the user, notes-61, 2026-10-06: "start thinking of \\
\texttt{library-book.scm} & THE BOOK OF THEOREMS AND DEFINITIONS, NO PROOFS (the user, \\
\texttt{library.scm} & theory management \\
\texttt{load.scm} & load the proof checker in dependency order \\
\texttt{macetes.scm} & macete mechanism \\
\texttt{major-theorems.scm} & the LaTeX list of MAJOR THEOREMS \\
\texttt{minimize.scm} & `minimize!', the least-measure tactic \\
\texttt{mint-attribution.scm} & \textit{(no header)} \\
\texttt{mp.scm} & `(mp)': close a goal that is an INSTANCE of a universal you already \\
\texttt{mutation-check.scm} & IS THE MUST-NOT-PROVE CORPUS ACTUALLY DETECTING ANYTHING? \\
\texttt{number-systems.scm} & number system constants and axioms \\
\texttt{operators.scm} & ONE table, keyed by head symbol \\
\texttt{page-audit-support.scm} & \textit{(no header)} \\
\texttt{page-audit.scm} & THE GATE ON THE PRINTED PAGE \\
\texttt{parser.scm} & VNB string expression parser \\
\texttt{parts.scm} & accessing the PARTS of a formula by SURFACE syntax \\
\texttt{preamble.scm} & rule-based drivers a human can edit (notes-45, batch 41, 2026-09-28), \\
\texttt{prep.scm} & (prep 'FUBA): why will FUBA not fire here, and what must be done \\
\texttt{presentation.scm} & THE RESOLUTION DIAL \\
\texttt{primitive-inferences.scm} & the primitive inference kernel \\
\texttt{proof-commands.scm} & higher-level proof commands (tactics) \\
\texttt{proof-debt.scm} & the warrant / proof-debt ledger \\
\texttt{proof-map.scm} & (proof-map): see the WHOLE proof tree, closed branches \\
\texttt{proof-reader.scm} & the HUMAN-READABLE proof sketch (``reader mode''), the \\
\texttt{proof-tex.scm} & render a completed proof as a LaTeX step-trace \\
\texttt{prop.scm} & `prop': decide a propositional entailment, then PROVE it \\
\texttt{proven-theorems.scm} & TOMBSTONE (2026-05-27) \\
\texttt{push-not.scm} & (push-not-h k): push a NOT inward through one quantifier or \\
\texttt{ra-diverge.scm} & \textit{(no header)} \\
\texttt{reference-html.scm} & the browser pages of the VNB reference, built by \\
\texttt{replay-audit.scm} & does a recorded proof SCRIPT replay to a grounded graph? \\
\texttt{rule-checkers-logic.scm} & the 42 `logic' kernel operations, CHECKED \\
\texttt{rule-checkers-oracle.scm} & inference checkers for the seven ORACLE rules \\
\texttt{rule-checkers-rewrite.scm} & checkers for the four `rewrite' kernel operations \\
\texttt{rule-checkers-schema.scm} & checkers for the thirteen `schema' kernel operations \\
\texttt{sequents.scm} & sequents over <wff> objects \\
\texttt{sketch.scm} & the structured-proof surface (sketch/step/obtain/qed-sketch), \\
\texttt{statement-types.scm} & THE IMPLIED TYPES OF A STATEMENT (the user's request, 2026-10-17) \\
\texttt{strapp-fast.scm} & \textit{(no header)} \\
\texttt{strapp-patch.scm} & MIT/GNU Scheme 12.1's string-append without its two 1024-slot \\
\texttt{structures.scm} & mathematical structure definitions \\
\texttt{suggest.scm} & assumption-pattern scanner for forward-move discovery \\
\texttt{tactic-uses-data.scm} & GENERATED by docs/gen-tactic-uses.py from \\
\texttt{tactics-help.scm} & a self-describing registry of the interactive tactics \\
\texttt{test-suite-negative.scm} & the MUST-NOT-PROVE corpus \\
\texttt{test-suite.scm} & VNB system regression and demonstration tests \\
\texttt{tex-infix.scm} & THE OPTIONAL INFIX LAYER for the TeX printer (printer side) \\
\texttt{tex-output.scm} & render VNB s-expression formulas as LaTeX math \\
\texttt{transport.scm} & TRANSPORTABLE FACTS \\
\texttt{vlet.scm} & (vlet (names...) FORMER): bind names to proof objects \\
\texttt{wff-english.scm} & render a wff as an English sentence \\
\texttt{wff.scm} & the <wff> record type (well-formed formula in a theory) \\
\texttt{witness-tactics.scm} & getting a witness out of a quantified hypothesis \\
\texttt{zero-it.scm} & (zero-it): move an equation goal P = Q to its normal form R = 0 \\
\bottomrule
\end{longtable}

\section{\texttt{structure-library/}}
\label{sec:src-structure-library}
The 322 files that install definitions, structures, vocabulary, and warranted supports.  One file per structure or cluster; the always-current per-structure listing (declaration, slots, parent, theorems) is the generated \texttt{reference/STRUCTURE-INDEX.md}.

\begin{longtable}{@{}p{0.30\textwidth}p{0.63\textwidth}@{}}
\toprule
\textbf{File} & \textbf{Description} \\
\midrule
\endhead
\texttt{abelian-group.scm} & ABELIAN-GROUP structure \\
\texttt{algnum-vocab.scm} & ALGEBRAIC ELEMENTS OF C (ALGNUM, algebraic-numbers bundle B7, 2026-10-26) \\
\texttt{analytic-log.scm} & ANALYTICITY AT A POINT (the notes' Def 2.11) AND THE \\
\texttt{antiderivative.scm} & Definition 4.6 (calculus.pdf), as VOCABULARY \\
\texttt{ascoli-arzela-statement.scm} & \textit{(no header)} \\
\texttt{ba-amp-vocab.scm} & Thayer, Banach Algebras and Applications, ch. 6 (the amplification in the proof of \\
\texttt{ba-bh2-vocab.scm} & wave of 2026-11-09/10 (agent BA-BH2), wired 2026-11-10 from scratchpad/babh2-1109/ba-bh2-vocab.scm \\
\texttt{ba-cx-vocab.scm} & wave of 2026-11-09/10 (agent BA-1), wired 2026-11-10 from scratchpad/ba1-1109/ba-cx-vocab.scm \\
\texttt{ba-gns-vocab.scm} & THE GNS CONSTRUCTION, the pre-Hilbert stage: BanachAlgebras.pdf Theorem 4.13 and \\
\texttt{ba-gns2-vocab.scm} & THE GNS CONSTRUCTION ON THE HILBERT COMPLETION and THE RADICAL: \\
\texttt{ba-neumann-vocab.scm} & wave of 2026-11-09/10 (agent BA-1), wired 2026-11-10 from scratchpad/ba1-1109/ba-neumann-vocab.scm \\
\texttt{ba-ops-vocab.scm} & wave of 2026-11-09/10 (agent BA-BH), wired 2026-11-10 from scratchpad/babh-1109/ba-ops-vocab.scm \\
\texttt{ba-optop-vocab.scm} & wave of 2026-11-09/10 (agent BA-OT1), wired 2026-11-10 from scratchpad/baot1-1109/ba-optop-vocab.scm \\
\texttt{ba-quotient-na-vocab.scm} & wave of 2026-11-09/10 (agent BA-1E), wired 2026-11-10 from scratchpad/ba1e-1109/ba-quotient-na-vocab.scm \\
\texttt{ba-quotient-vocab.scm} & wave of 2026-11-09/10 (agent BA-1C), wired 2026-11-10 from scratchpad/ba1c-1109/ba-quotient-vocab.scm \\
\texttt{ba-rep-vocab.scm} & POSITIVE FUNCTIONALS, STATES AND REPRESENTATIONS of a unital involutive \\
\texttt{ba-trace-dual-vocab.scm} & Thayer, Banach Algebras and Applications, chapter 5 s.6-s.8 (spaces of \\
\texttt{ba-unbdd-vocab.scm} & wave of 2026-11-09/10 (agent BA-G1), wired 2026-11-10 from scratchpad/bag1-1109/ba-unbdd-vocab.scm \\
\texttt{ba-vn-vocab.scm} & Thayer, Banach Algebras and Applications, chapter 6 (density theorems): the \\
\texttt{ba-vocab.scm} & wave of 2026-11-09/10 (agent BA-1), wired 2026-11-10 from scratchpad/ba1-1109/ba-vocab.scm \\
\texttt{ba2-vocab.scm} & wave of 2026-11-09/10 (agent BA-2), wired 2026-11-10 from scratchpad/ba2-1109/ba2-vocab.scm \\
\texttt{ba2b-gen-vocab.scm} & wave of 2026-11-09/10 (agent BA-2B), wired 2026-11-10 from scratchpad/ba2b-1109/ba2b-gen-vocab.scm \\
\texttt{baire-category.scm} & interior/closure, nowhere-dense, meager/nonmeager, and \\
\texttt{basw-vocab.scm} & wave of 2026-11-09/10 (agent BA-SW), wired 2026-11-10 from scratchpad/basw-1109/basw-vocab.scm \\
\texttt{basw2-vocab.scm} & wave of 2026-11-09/10 (agent BA-SW2), wired 2026-11-10 from scratchpad/basw2-1109/basw2-vocab.scm \\
\texttt{basw3-vocab.scm} & wave of 2026-11-09/10 (agent BA-SW3), wired 2026-11-10 from scratchpad/basw3-1109/basw3-vocab.scm \\
\texttt{bidual.scm} & \textit{(no header)} \\
\texttt{bijection-derived.scm} & the three BIJECTION projection lemmas, PROVEN \\
\texttt{bijection.scm} & bijections between classes, and INVERSE-BIJ via CHOICE \\
\texttt{binomial.scm} & the Binomial Theorem for commutative rings, over the \\
\texttt{block-family-combinatorial.scm} & \textit{(no header)} \\
\texttt{bounded-metric.scm} & the bounded metric  rho = d/(1+d)  of a metric space, \\
\texttt{c-metric-space.scm} & C-METRIC-SPACE, a set carrying a COUNTABLE FAMILY OF \\
\texttt{ca-ac-vocab.scm} & UNIFORM CONVERGENCE ON COMPACT SUBSETS of an open set of CC \\
\texttt{ca-angle-vocab.scm} & WAVE 2 (2026-10-27..30), wired 2026-10-30 from scratchpad/ch3e/ch3e-angle.scm (vocab) \\
\texttt{ca-homotopy-vocab.scm} & \textit{(no header)} \\
\texttt{ca-hopf-vocab.scm} & THE HOPF FIBRATION, complex-analysis.pdf s. 1.6.3 ((13)-(16), Prop. 1.22, \\
\texttt{ca-lft-vocab.scm} & THE EXTENDED PLANE AND LINEAR FRACTIONAL TRANSFORMATIONS, \\
\texttt{ca-line-vocab.scm} & WAVE 2 (2026-10-27..30), wired 2026-10-30 from scratchpad/ch3e/ch3f-line.scm (vocab), scratchpad/ch3e/ch3f-refl.scm (vocab) \\
\texttt{ca1-vocab.scm} & THE COMPLEX NUMBERS AS 2-BY-2 REAL MATRICES (complex-analysis.pdf, \\
\texttt{ca5-function-spaces.scm} & THE VOCABULARY OF CHAPTER 5 (FUNCTION SPACES) of the \\
\texttt{ca51-vocab.scm} & local compactness and exhaustions by compact sets (CA-51, 2026-10-26) \\
\texttt{calc3-vocab.scm} & vocabulary for calculus.pdf chapter 3 (metric topology): \\
\texttt{calint-vocab.scm} & WAVE 2 (2026-10-27..30), wired 2026-10-30 from scratchpad/calint/calint-vocab.scm \\
\texttt{cardinality.scm} & CARD, the cardinality of a set, DEFINED \\
\texttt{catools-gal-vocab.scm} & divisibility in an integral domain \\
\texttt{cauchy-subsequence.scm} & \textit{(no header)} \\
\texttt{cc-coords.scm} & CC and cartesian(RR, RR), as VOCABULARY \\
\texttt{cc-elementary.scm} & THE ELEMENTARY FUNCTIONS exp, cos, sin ON CC, DEFINED BY \\
\texttt{cc-pi.scm} & THE CONSTANT pi, DEFINED FROM THE KERNEL OF exp (the notes 2.3.3) \\
\texttt{cc-power-series.scm} & THE VOCABULARY OF COMPLEX POWER SERIES (the notes 2.2) \\
\texttt{ccr-ch3-vocab.scm} & wave of 2026-11-09/10 (agent CCR-1109-B), wired 2026-11-10 from scratchpad/ccr1109b/ccr-ch3-vocab.scm \\
\texttt{ccr-ch3b-vocab.scm} & wave of 2026-11-09/10 (agent CCR-1109-B2), wired 2026-11-10 from scratchpad/ccr1109b2/ccr-ch3b-vocab.scm \\
\texttt{ccr-ch3c-vocab.scm} & wave of 2026-11-09/10 (agent CCR-1109-B3), wired 2026-11-10 from scratchpad/ccr1109b3/ccr-ch3c-vocab.scm \\
\texttt{ccr1-vocab.scm} & ULTRAMETRIC SPACES (CCR.pdf, chapter 1, Definition 1.1 and the objects of \\
\texttt{ccr1109a2-vocab.scm} & wave of 2026-11-09/10 (agent CCR-1109-A2), wired 2026-11-10 from scratchpad/ccr1109a2/ccr1109a2-vocab.scm \\
\texttt{ccr1109a3-laurent-vocab.scm} & wave of 2026-11-09/10 (agent CCR-1109-A3), wired 2026-11-10 from scratchpad/ccr1109a3/ccr1109a3-laurent-vocab.scm \\
\texttt{ccr1109a4-laurent-lim-vocab.scm} & wave of 2026-11-09/10 (agent CCR-1109-A4), wired 2026-11-10 from scratchpad/ccr1109a4/ccr1109a4-laurent-lim-vocab.scm \\
\texttt{ccr1109a5-vocab.scm} & wave of 2026-11-09/10 (agent CCR-1109-A5), wired 2026-11-10 from scratchpad/ccr1109a5/ccr1109a5-vocab.scm \\
\texttt{ccr1109a6-vocab.scm} & wave of 2026-11-09/10 (agent CCR-1109-A6), wired 2026-11-10 from scratchpad/ccr1109a6/ccr1109a6-vocab.scm \\
\texttt{ccr2-vocab.scm} & NON-ARCHIMEDEAN VALUED FIELDS AND THEIR DOMAINS (CCR.pdf chapters 1-2 and 4: \\
\texttt{ccr3-vocab.scm} & MODULES IN A LOCAL NV FIELD (CCR.pdf Definition 4.11) \\
\texttt{ccr4-vocab.scm} & ADDITIVE CHARACTERS OF A NORMED FIELD AND THE DUALITY MAP (CCR.pdf chapter 5, Definitions \\
\texttt{ccr4b-vocab.scm} & wave of 2026-11-09/10 (agent CCR-1109-C3), wired 2026-11-10 from scratchpad/ccr1109c3/ccr4b-vocab.scm \\
\texttt{ccr4c-subnf-vocab.scm} & wave of 2026-11-09/10 (agent CCR-1109-C4), wired 2026-11-10 from scratchpad/ccr1109c4/ccr4c-subnf-vocab.scm \\
\texttt{ccr4d-ext-vocab.scm} & wave of 2026-11-09/10 (agent CCR-1109-C5), wired 2026-11-10 from scratchpad/ccr1109c5/ccr4d-ext-vocab.scm \\
\texttt{ccr4f-vocab.scm} & THE DIGIT MAP AND THE COORDINATE MAP AS FUNCTIONS (CCR-1111 finding F1) \\
\texttt{ccr4x-ext-vocab.scm} & wave of 2026-11-09/10 (agent CCR-1109-C2), wired 2026-11-10 from scratchpad/ccr1109c2/ccr4x-ext-vocab.scm \\
\texttt{ccr4x-vocab.scm} & wave of 2026-11-09/10 (agent CCR-1109-C), wired 2026-11-10 from scratchpad/ccr1109c/ccr4x-vocab.scm \\
\texttt{ccr5-vocab.scm} & wave of 2026-11-09/10 (agent CCR-1109-D), wired 2026-11-10 from scratchpad/ccr1109d/ccr5-vocab.scm \\
\texttt{ccr6e-vocab.scm} & wave of 2026-11-09/10 (agent CCR-1109-D3), wired 2026-11-10 from scratchpad/ccr1109d3/ccr6e-vocab.scm \\
\texttt{cg1-vocab.scm} & vocabulary for the compact-groups notes, chapter 5 section 2 (CG-1) \\
\texttt{cg2-vocab.scm} & set products and set powers in a group (compact-groups notes, Prop 5.8: \\
\texttt{cg4-vocab.scm} & compact-groups notes, sections 1 and 5 (CG-4 session) \\
\texttt{ch3a-vocab.scm} & automorphisms of an open subset of CC (CH3-A, 2026-10-26) \\
\texttt{ch4a-vocab.scm} & WAVE 2 (2026-10-27..30), wired 2026-10-30 from scratchpad/ch4a/ch4a-vocab.scm \\
\texttt{ch4b-vocab.scm} & WAVE 2 (2026-10-27..30), wired 2026-10-30 from scratchpad/ch4b/ch4b-vocab.scm \\
\texttt{ch4d-vocab.scm} & WAVE 2 (2026-10-27..30), wired 2026-10-30 from scratchpad/ch4d/ch4d-vocab.scm \\
\texttt{cnvs-ba-vocab.scm} & wave of 2026-11-09/10 (agent CNVS), wired 2026-11-10 from scratchpad/cnvs-1109/cnvs-ba-vocab.scm \\
\texttt{cnvs-comp-vocab.scm} & wave of 2026-11-09/10 (agent CNVS2), wired 2026-11-10 from scratchpad/cnvs2-1109/cnvs2-comp-vocab.scm \\
\texttt{cnvs-vocab.scm} & wave of 2026-11-09/10 (agent CNVS), wired 2026-11-10 from scratchpad/cnvs-1109/cnvs-vocab.scm \\
\texttt{comm-ring-simplify.scm} & commutative-ring identity decision procedure \\
\texttt{commutative-ring.scm} & COMMUTATIVE-RING: a ring whose MUL is commutative \\
\texttt{compactness.scm} & compactness of a metric space and the four-way \\
\texttt{complex-inner-product.scm} & COMPLEX-INNER-PRODUCT-SPACE: a vector space over \\
\texttt{complex-measure.scm} & WAVE 2 (2026-10-27..30), wired 2026-10-30 from scratchpad/mf/mf-cvocab.scm \\
\texttt{complex.scm} & CC as a metric space \\
\texttt{compose-typing.scm} & RAN (range) + mindless type-closure lemmas for \\
\texttt{compose.scm} & COMPOSE: composition of functions, f o g \\
\texttt{convex-functional.scm} & \textit{(no header)} \\
\texttt{cvx-vocab.scm} & convex combinations and the convex hull in a real normed vector space \\
\texttt{definitional-reclass.scm} & reclassify hand-written DEFINING iffs \\
\texttt{derivative.scm} & the Caratheodory derivative, as VOCABULARY \\
\texttt{determinant.scm} & the determinant of a square matrix over a ring, by \\
\texttt{diff-on.scm} & THE DERIVATIVE ON AN OPEN SET, OVER A NORMED FIELD \\
\texttt{doubling-functions.scm} & WAVE 2 (2026-10-27..30), wired 2026-10-30 from scratchpad/mf/mf-mivocab.scm \\
\texttt{dual-space.scm} & THE DUAL of a real normed vector space, as an OBJECT \\
\texttt{elementary-matrix.scm} & matrix units and elementary column operations \\
\texttt{euclidean-ring.scm} & EUCLIDEAN-RING: an integral domain admitting a \\
\texttt{extended-arith.scm} & the rest of the arithmetic and order vocabulary of \\
\texttt{extended-reals-pos.scm} & RR-POS-STAR = [0, +inf], the nonnegative extended reals \\
\texttt{extended-reals-updown.scm} & \textit{(no header)} \\
\texttt{extended-reals.scm} & the extended real line RR-STAR = [-inf, +inf] \\
\texttt{extended-sum.scm} & \textit{(no header)} \\
\texttt{extreme-value.scm} & the closed interval [a,b] and the Extreme Value Theorem \\
\texttt{field.scm} & FIELD: a commutative ring in which every nonzero element \\
\texttt{finite-dimensional.scm} & finite-dimensional vector spaces, Zorn-free \\
\texttt{finprod.scm} & FINPROD / PROD-RING: finite PRODUCT of a function over a \\
\texttt{finsum.scm} & FINSUM: summation of an abelian-group-valued function \\
\texttt{flat-regular.scm} & PUBLIC-KEY SPACES WHOSE REGULAR CONTENTS ARE FLAT: the one property of the \\
\texttt{frechet-banach-steinhaus.scm} & BOUNDED SETS and EQUICONTINUOUS FAMILIES of linear \\
\texttt{frechet-metric.scm} & the METRIC of a seminormed space whose seminorm family is \\
\texttt{frechet-open-mapping.scm} & Frechet spaces, continuous linear maps, and the \\
\texttt{functor-invariance.scm} & a construction that cannot see the functor \\
\texttt{functoriality.scm} & every view-as is a FUNCTOR, and here is the proof \\
\texttt{gal-ch-vocab.scm} & vocabulary for the general Cayley-Hamilton theorem (GALOIS B6, 2026-10-26) \\
\texttt{gal-char-vocab.scm} & prime natural numbers and the characteristic of a field (galois-theory notes, \\
\texttt{gal-charpoly-vocab.scm} & the companion matrix and the characteristic polynomial (galois-theory notes, \\
\texttt{gal-gcd-vocab.scm} & a greatest common divisor of a set of ring elements (galois-theory notes, \\
\texttt{gal-op-vocab.scm} & a linear operator T on a module E over F as an F[x]-module: the iterates T\textasciicircum{}k, the \\
\texttt{gal-poly-vocab.scm} & evaluation of a formal polynomial and the polynomial x - c \\
\texttt{gal-ufd-vocab.scm} & associates and unique factorization domains (galois-theory notes, section 1) \\
\texttt{gamma-continuation-vocab.scm} & WAVE 2 (2026-10-27..30), wired 2026-10-30 from scratchpad/ch3g/ch3g-gamma.scm (vocab) \\
\texttt{gamma-function.scm} & THE IMPROPER INTEGRAL ON (0, +inf) AND THE GAMMA FUNCTION \\
\texttt{gn-vocab.scm} & BANACH ALGEBRAS: the natural order of an involutive algebra (Def. 1.34, (15)) and the \\
\texttt{group-cosets-vocab.scm} & \textit{(no header)} \\
\texttt{group-zpow-vocab.scm} & integer powers in a group (ALGNUM, 2026-10-27) \\
\texttt{group.scm} & GROUP structure \\
\texttt{hahn-banach-sublinear.scm} & \textit{(no header)} \\
\texttt{hausdorff-vocab.scm} & \textit{(no header)} \\
\texttt{hc-vocab.scm} & wave of 2026-11-09/10 (agent HAAR-C), wired 2026-11-10 from scratchpad/haarc-1109/hc-vocab.scm \\
\texttt{hc2-vocab.scm} & wave of 2026-11-09/10 (agent HAAR-C2), wired 2026-11-10 from scratchpad/haarc2-1109/hc2-vocab.scm \\
\texttt{hc3-vocab.scm} & wave of 2026-11-09/10 (agent HAAR-C3), wired 2026-11-10 from scratchpad/haarc3-1109/hc3-vocab.scm \\
\texttt{hilbert-completion-vocab.scm} & \textit{(no header)} \\
\texttt{hilbert-onb-vocab.scm} & wave of 2026-11-09/10 (agent BA-ONB), wired 2026-11-10 from scratchpad/baonb-1109/hilbert-onb-vocab.scm \\
\texttt{hilbert-space.scm} & HILBERT SPACES: the metric space of a complex inner \\
\texttt{id-fun.scm} & ID-FUN: the identity map of a set \\
\texttt{ideal.scm} & ideals, principal ideals, and principal-ideal domains over a \\
\texttt{ineq-oracle.scm} & the (ineq) decision procedure: close a linear-inequality \\
\texttt{inf-subsets.scm} & the class of infinite subsets of a given set \\
\texttt{infinite-products-vocab.scm} & WAVE 2 (2026-10-27..30), wired 2026-10-30 from scratchpad/ch6/ch6-prod.scm (vocab) \\
\texttt{injection.scm} & INJECTION class, IMAGE operator, and the \\
\texttt{integral-domain-laws.scm} & integral-domain-cancel-zero, PROVEN from the \\
\texttt{integral-domain.scm} & INTEGRAL-DOMAIN: a commutative ring with ONE /= ZERO \\
\texttt{integral.scm} & \textit{(no header)} \\
\texttt{intersection-of.scm} & INTERSECTION-OF(c), the intersection of the members \\
\texttt{interval-calculus.scm} & THE VOCABULARY OF THE CALCULUS ON AN INTERVAL \\
\texttt{interval-extremum.scm} & LOCAL MAXIMUM AND LOCAL MINIMUM OF A FUNCTION ON \\
\texttt{interval-taylor.scm} & THE VOCABULARY OF TAYLOR'S FORMULA ON AN INTERVAL \\
\texttt{ipf-vocab.scm} & INNER PRODUCT FORMS ON A COMPLEX MODULE, the module carrying such a form, and the \\
\texttt{is-set-laws.scm} & \textit{(no header)} \\
\texttt{isom-vocab.scm} & vocabulary for calculus.pdf ch. 3, sections 6 and 8: uniform and \\
\texttt{isom2-rrvs-vocab.scm} & the real numbers as a FIELD instance, the vector space of \\
\texttt{isom2-vocab.scm} & the group structure of the isometry group (calculus.pdf ch. 3 \\
\texttt{km-vocab.scm} & wave of 2026-11-09/10 (agent BA-KM), wired 2026-11-10 from scratchpad/bakm-1109/km-vocab.scm \\
\texttt{l1-dual.scm} & WAVE 2 (2026-10-27..30), wired 2026-10-30 from scratchpad/md/me-du-vocab.scm \\
\texttt{l2-families-vocab.scm} & wave of 2026-11-09/10 (agent BA-ONB2), wired 2026-11-10 from scratchpad/baonb2-1109/bo2-l2-vocab.scm \\
\texttt{l2-integral-operator.scm} & wave of 2026-11-09/10 (agent TH3-L2C), wired 2026-11-10 from scratchpad/th3-1109/l2c/kc-vocab.scm \\
\texttt{l2c-vocab.scm} & wave of 2026-11-09/10 (agent TH3-L2), wired 2026-11-10 from scratchpad/th3-1109/l2/l2c-vocab.scm \\
\texttt{la-det-vocab.scm} & triangular matrices and the formal adjoint (LA-DET, 2026-10-18) \\
\texttt{la-elem-vocab.scm} & elementary column matrices as a class, the subgroup they \\
\texttt{la-z-vocab.scm} & subdeterminants, equivalence by elementary operations and the \\
\texttt{la-z2-vocab.scm} & the principal sequence of a subgroup (linear algebra, the \\
\texttt{la-z3-vocab.scm} & quotient modules, the ring as a module over itself, the cyclic \\
\texttt{la-z4-vocab.scm} & the free module Z\textasciicircum{}n, the image of a sequence under a map, the \\
\texttt{la-z5-vocab.scm} & the norm N\_E(F) of a submodule (linear algebra, the definition \\
\texttt{la-z6-vocab.scm} & polynomials in n indeterminates and the permutation of the indeterminates \\
\texttt{lebesgue-translate.scm} & the TRANSLATE of a set of reals \\
\texttt{line-coefficient.scm} & \textit{(no header)} \\
\texttt{linear-arith.scm} & standalone linear-arithmetic decision core \\
\texttt{linear-functional.scm} & bounded linear functionals and the dual norm on a \\
\texttt{list-recursion.scm} & CONS, and the recursive characterization of LENGTH \\
\texttt{little-o.scm} & calculus.pdf Section 2.2, the calculus of o/O notation, in a \\
\texttt{local-measure.scm} & WAVE 2 (2026-10-27..30), wired 2026-10-30 from scratchpad/mf/mf-mgvocab.scm \\
\texttt{lp-inequalities.scm} & POWERS IN [0,+inf] AND THE p-TH POWER INTEGRAL I\_p \\
\texttt{lp-seminorm.scm} & THE L\textasciicircum{}p SEMINORM OF A REAL-VALUED FUNCTION \\
\texttt{m-riesz-vocab.scm} & ORDERED VECTOR SPACES: BanachAlgebras.pdf section 4.6, before Proposition 4.32 \\
\texttt{mat-equiv.scm} & the matrix EQUIVALENCE relation \textasciitilde{} (algebraic-numbers.pdf \\
\texttt{matrix.scm} & \textit{(no header)} \\
\texttt{mb-vocab.scm} & vocabulary for the M-B wave (construction-of-measures.pdf 2.9, 3.2-3.4, \\
\texttt{mb12-vocab.scm} & vocabulary for construction-of-measures.pdf Appendix A.3: the \\
\texttt{mb35-souslin-vocab.scm} & vocabulary for construction of measures, chapter 5 section 1.3-1.4: sieves and \\
\texttt{mb35-vocab.scm} & vocabulary for construction of measures, chapter 5 section 2 (M-B35, 2026-10-26) \\
\texttt{mb4-vocab.scm} & vocabulary for construction of measures, chapter 5 (M-B4, 2026-10-26) \\
\texttt{mb67-vocab.scm} & vocabulary for construction of measures, chapter 5 sections 1.4-1.6: the Baire space N\textasciicircum{}N \\
\texttt{measurable-space.scm} & MEASURABLE-SPACE and MEASURE-SPACE as declared \\
\texttt{measure-ch1-vocab.scm} & ATOMS OF AN ALGEBRA, FINITE UNIONS, RECTANGLES OF R\textasciicircum{}n, THE PULLBACK SIGMA-ALGEBRA AND \\
\texttt{measure-step.scm} & WAVE 2 (2026-10-27..30), wired 2026-10-30 from scratchpad/md/md-vocab.scm \\
\texttt{measure.scm} & OUTER MEASURES, CARATHEODORY MEASURABILITY, MEASURES \\
\texttt{meta-core.scm} & THE CORE OF K2: well-formedness, sequents, the core rule relations, \\
\texttt{meta-cr.scm} & CONTEXTUAL REWRITING (CR) over the encoded core: the one-step rewrite relation \\
\texttt{meta-def.scm} & K2 PHASE 5: the clauses of the kernel's definedness certificate pi--defined? \\
\texttt{meta-ind.scm} & K2 PHASE 4: the kernel relations of the induction rules, the set-formation rules \\
\texttt{meta-nbg.scm} & K2, agent METACORE (2026-10-29): the kernel forms of the ten C+ rows of the K2 ledger \\
\texttt{meta-syntax.scm} & THE SYNTAX OF THE KERNEL AS A VNB OBJECT: locally nameless expressions, \\
\texttt{metric-completeness.scm} & Cauchy sequences, convergence, completeness \\
\texttt{metric-completion.scm} & \textit{(no header)} \\
\texttt{metric-connected.scm} & CONNECTEDNESS OF A SUBSET OF A METRIC SPACE \\
\texttt{metric-continuity.scm} & continuous maps between metric spaces \\
\texttt{metric-laws.scm} & the five metric laws, PROVEN from the definition \\
\texttt{metric-open-sets.scm} & \textit{(no header)} \\
\texttt{metric-outer-measure.scm} & the distance between sets, METRIC OUTER MEASURES and \\
\texttt{metric-space.scm} & METRIC-SPACE structure \\
\texttt{metric-subspace.scm} & the METRIC SUBSPACE of a metric space, and the \\
\texttt{metric-topology.scm} & balls, r-nets, total-boundedness \\
\texttt{minkowski-functional.scm} & \textit{(no header)} \\
\texttt{mod-seq.scm} & matrices acting on sequences of module elements \\
\texttt{module.scm} & MODULE over a ring, with the scalar action ACT \\
\texttt{monoid-power.scm} & MPOW: the natural-number power x\textasciicircum{}n in a monoid \\
\texttt{monoid.scm} & MONOID and COMM-MONOID structures \\
\texttt{monotone-class.scm} & MONOTONE CLASSES \\
\texttt{mvp-vocab.scm} & WAVE 2 (2026-10-27..30), wired 2026-10-30 from scratchpad/calint2/mvp-vocab.scm \\
\texttt{nets.scm} & \textit{(no header)} \\
\texttt{nn-arith.scm} & the recursion equations for + and * on NN \\
\texttt{nn-minus.scm} & truncated natural subtraction (monus), as VOCABULARY \\
\texttt{normed-ag-metric.scm} & the metric space underlying a normed abelian group \\
\texttt{normed-ag.scm} & NORMED-AG structure (normed abelian group) \\
\texttt{normed-field-metric.scm} & the metric space underlying a normed field \\
\texttt{normed-field.scm} & NORMED-FIELD: a commutative ring with multiplicative \\
\texttt{normed-vector-space.scm} & NORMED-VECTOR-SPACE: a real normed vector space \\
\texttt{nsl.scm} & THE NEEDHAM-SCHROEDER-LOWE PROTOCOL ON THE USER'S STRAND MODEL: names, public keys, \\
\texttt{nso.scm} & THE ORIGINAL NEEDHAM-SCHROEDER PROTOCOL (Thayer, Herzog, Guttman, section 5.3), ON THE \\
\texttt{nth-deriv.scm} & THE ITERATED DERIVATIVE ON AN OPEN SET, OVER A NORMED FIELD \\
\texttt{numeric-instances.scm} & \textit{(no header)} \\
\texttt{nvlocal-vocab.scm} & WAVE 2 (2026-10-27..30), wired 2026-10-30 from scratchpad/ccrqp/nvlocal-vocab.scm \\
\texttt{nvs-metric.scm} & the metric induced by the norm of a real normed vector \\
\texttt{one-sided-derivative.scm} & THE ONE-SIDED DERIVATIVES OF THE USER'S NOTES, \\
\texttt{operation-properties.scm} & named properties of operations \\
\texttt{or-missing-guarantee-instance.scm} & \textit{(no header)} \\
\texttt{order-lemmas.scm} & the finite order calculus of RR \\
\texttt{order-predicates.scm} & derived order predicates on the kernel <= \\
\texttt{order-zorn.scm} & abstract partial orders and ZORN'S LEMMA (statement only) \\
\texttt{ordinals.scm} & ordinal numbers: axioms, ordering, induction, recursion \\
\texttt{otway-rees.scm} & THE OTWAY-REES PROTOCOL ON THE USER'S STRAND MODEL (Thayer, Herzog, Guttman, \\
\texttt{outer-measure-construction.scm} & METHOD I OUTER MEASURES \\
\texttt{padic-vocab.scm} & THE p-ADIC ABSOLUTE VALUE ON QQ (CCR.pdf section 1.2) \\
\texttt{path-integral.scm} & THE VOCABULARY OF THE INTEGRAL ALONG A PATH \\
\texttt{poly-degree.scm} & DEG, LEADCOEF and MONOMIAL for the monoid algebra's \\
\texttt{poly-fn-vocab.scm} & POLYNOMIAL FUNCTIONS (CA-MONTEL ``POLY'', 2026-10-26) \\
\texttt{polynomial.scm} & the monoid-algebra construction A[M], and polynomials as \\
\texttt{positive-functional.scm} & \textit{(no header)} \\
\texttt{power-series.scm} & real power series  Sum\_\{n>=0\} coef(n) x\textasciicircum{}n  in RR \\
\texttt{prod-of-sums.scm} & \textit{(no header)} \\
\texttt{product-measure.scm} & PRODUCT SIGMA-ALGEBRAS, SECTIONS, DYNKIN SYSTEMS \\
\texttt{product-metric.scm} & the countable product of metric spaces, as a metric \\
\texttt{product2-vocab.scm} & THE BINARY PRODUCT OF TWO METRIC SPACES (definitions only) \\
\texttt{pruefer-vocab.scm} & WAVE 2 (2026-10-27..30), wired 2026-10-30 from scratchpad/ccrpr/pruefer-vocab.scm \\
\texttt{pseudometric.scm} & PSEUDOMETRIC-SPACE, the topology GENERATED by a countable \\
\texttt{push-forward.scm} & \textit{(no header)} \\
\texttt{qag-vocab.scm} & WAVE 2 (2026-10-27..30), wired 2026-10-30 from scratchpad/ccrpr/qag-vocab.scm \\
\texttt{qp-vocab.scm} & WAVE 2 (2026-10-27..30), wired 2026-10-30 from scratchpad/ccrqp/qp-vocab.scm \\
\texttt{qq-fractions.scm} & QQ is the field of fractions of the integers \\
\texttt{rabasics-vocab.scm} & vocabulary for the RA-BASICS wave (calculus.pdf 1.13, 3.7 \\
\texttt{rb2-vocab.scm} & vocabulary for calculus.pdf 2.27 and for path-connectedness \\
\texttt{real-integral.scm} & THE INTEGRAL OF A REAL-VALUED FUNCTION: positive and \\
\texttt{real-powers.scm} & the real power RPOW and the square root SQRT \\
\texttt{reduce.scm} & REDUCE: left-fold of a binary op over a finite indexed \\
\texttt{references.scm} & \textit{(no header)} \\
\texttt{regulated-primitive.scm} & COUNTABLE SETS, ONE-SIDED LIMITS WITHIN A SET, \\
\texttt{regulated.scm} & docs/calculus.pdf CHAPTER 4 SECTION 1: the vocabulary of \\
\texttt{riesz-markov-model.scm} & \textit{(no header)} \\
\texttt{ring-pow-vocab.scm} & wave of 2026-11-09/10 (agent CCR-1109-E), wired 2026-11-10 from scratchpad/ccr1109e/ring-pow-vocab.scm \\
\texttt{ring-power.scm} & RING-POWER: the natural-number power x\textasciicircum{}n in a commutative ring \\
\texttt{ring-simplify.scm} & polynomial-normal-form ring identity decision procedure \\
\texttt{ring.scm} & RING, RING-PROD, ZERO-RING \\
\texttt{ringoid.scm} & a ring with a distinguished two-sided ideal \\
\texttt{rn-vocab.scm} & E = R\textasciicircum{}n WITH THE SUP NORM, calculus.pdf ch. 2 s. 8 (32), (36)-(37), (44) \\
\texttt{road-juxtaposition.scm} & THE JUXTAPOSITION OF ROADS, THE TRIANGULAR ROAD AND THE LENGTH \\
\texttt{rr-ineq.scm} & PROTOTYPE of the rr-ineq applier (plan item 2, the "dumb \\
\texttt{scalar-inequalities.scm} & a library of important elementary inequalities \\
\texttt{segments-triangles.scm} & THE VOCABULARY OF CAUCHY-GOURSAT (complex-analysis.pdf \\
\texttt{semigroup.scm} & SEMIGROUP structure \\
\texttt{seminorm-extension.scm} & \textit{(no header)} \\
\texttt{seminorm-hahn-banach.scm} & seminorms on a K-vector space and the \\
\texttt{separable.scm} & separability of a metric space (a dense sequence), the \\
\texttt{seq-compact-product.scm} & \textit{(no header)} \\
\texttt{sequences.scm} & PROD-ORD, SUM, SUM-AG: products and sums over initial NN-segments \\
\texttt{set-basics.scm} & the NBG facts about SET that the kernel rules do not give \\
\texttt{set-vocabulary.scm} & two set constructors that the rest of the library is \\
\texttt{setoid.scm} & SETOID = (X, REL): a set with an equivalence relation, \\
\texttt{sigma-algebra.scm} & ALGEBRAS and SIGMA-ALGEBRAS of sets \\
\texttt{signed-measure.scm} & WAVE 2 (2026-10-27..30), wired 2026-10-30 from scratchpad/mf/mf-vocab.scm \\
\texttt{simple-integral.scm} & THE INTEGRAL OF A [0,+inf]-VALUED FUNCTION, DEFINED \\
\texttt{sip-complex-vocab.scm} & THE COMPLEX STRUCTURE J on a semi-inner-product space (SIP-1, \\
\texttt{sip-hilbert-vocab.scm} & REAL HILBERT SPACES: the norm and the metric of a \\
\texttt{sip-l2-vocab.scm} & EQUALITY ALMOST EVERYWHERE as a setoid (SIP-1, 2026-10-18) \\
\texttt{sip-vocab.scm} & SEMI-INNER-PRODUCT SPACES (SIP-1, 2026-10-18) \\
\texttt{sos-arith.scm} & exact-rational LP feasibility with witness extraction \\
\texttt{sos-oracle.scm} & the (sos) sum-of-squares oracle: close a nonstrict \\
\texttt{sqn.scm} & SQN(a): the sequences in a, i.e. FUN(NN, a) \\
\texttt{step-maps.scm} & WAVE 2 (2026-10-27..30), wired 2026-10-30 from scratchpad/mf/mf-mcvocab.scm \\
\texttt{strand-adversary.scm} & THE ADVERSARY: its traces, the infiltrated strand space \\
\texttt{strand-honesty.scm} & HONEST IDEALS (Thayer, Herzog, Guttman, \textasciitilde{}/docs/jcs\_strand\_spaces.pdf, Definition 6.10), \\
\texttt{strand-ideals.scm} & IDEALS OF THE MESSAGE ALGEBRA (Thayer, Herzog, Guttman, "Strand spaces: \\
\texttt{strand-messages.scm} & THE MESSAGES OF STRAND SPACES, ON THE USER'S MODEL, WITH ENCRYPTION \\
\texttt{strand-spaces.scm} & STRAND SPACES, NODES, BUNDLES ON THE USER'S MODEL \\
\texttt{subsequence-capture.scm} & \textit{(no header)} \\
\texttt{subtype-laws.scm} & the trivial subtype-subsumption laws, PROVEN \\
\texttt{summability.scm} & unconditional summability of a normed-AG-valued function \\
\texttt{tb11-vocab.scm} & vocabulary for topological-vector-spaces Appendix A.3 (TB-11): \\
\texttt{tg-vocab.scm} & topological groups (topological-vector-space.pdf Appendix A.3, TB-11) \\
\texttt{th3-vocab.scm} & wave of 2026-11-09/10 (agent TH3-A), wired 2026-11-10 from scratchpad/th3-1109/th3-vocab.scm \\
\texttt{th3a2-vocab.scm} & wave of 2026-11-09/10 (agent TH3-A2), wired 2026-11-10 from scratchpad/th3-1109/a2/th3a2-vocab.scm \\
\texttt{top-closed.scm} & \textit{(no header)} \\
\texttt{top-continuity.scm} & COMPACTNESS of a topological space, and the SUBSPACE \\
\texttt{top-product.scm} & \textit{(no header)} \\
\texttt{top-space.scm} & TOPOLOGICAL SPACES, and the functor Met -> Top \\
\texttt{transpose.scm} & the TRANSPOSE of a matrix (batch 32, 2026-09-25 \\
\texttt{trunc-metric.scm} & the TRUNCATED METRIC  d'(x,y) = min(1, d(x,y))  of a \\
\texttt{tvs-product-line.scm} & \textit{(no header)} \\
\texttt{tvs-vocab.scm} & \textit{(no header)} \\
\texttt{user-additions.scm} & structures saved via Build Structure \\
\texttt{views.scm} & view-as declarations between structures \\
\texttt{weak-star-separable.scm} & \textit{(no header)} \\
\texttt{weak-star.scm} & \textit{(no header)} \\
\texttt{winding-number.scm} & THE WINDING NUMBER AND THE CIRCLE AS A ROAD \\
\texttt{wot-compact-vocab.scm} & agent BA-1117-TRACE (round 2), 2026-11-17 \\
\texttt{zero-order-vocab.scm} & WAVE 2 (2026-10-27..30), wired 2026-10-30 from scratchpad/ch6/zero-order-vocab.scm \\
\texttt{zz-action.scm} & ZZ-ACT: the integer action on an abelian group \\
\texttt{zz-arith.scm} & ZZ is GENERATED by NN \\
\texttt{zz-divisibility.scm} & divisibility, Bezout combinations, gcd and coprimality \\
\texttt{zz-nt-vocab.scm} & \textit{(no header)} \\
\bottomrule
\end{longtable}

\section{\texttt{theorem-library/} and \texttt{calculus/}}
\label{sec:src-libraries}
\texttt{theorem-library/} holds 1112 proof scripts, each reaching \texttt{(qed \ldots)} and gated into the library load; by convention \texttt{X-proof.scm} proves \texttt{X}.  \texttt{calculus/} holds 21 probes and stress tests that are \emph{not} in \texttt{load.scm}.  Rather than enumerate them here, the theorems they prove are indexed, with statements and per-file links, in the generated browser references \texttt{reference/THEOREMS.md} (flat) and \texttt{reference/BY-TOPIC.md} (grouped like a textbook table of contents).

\section{Emacs interface (\texttt{emacs/}) and generated docs}
\label{sec:src-emacs}
The workspace front end lives in \texttt{emacs/}:
\begin{itemize}
\item \texttt{core-replay-stub.scm}
\item \texttt{vnb-cmd-gen.el}
\item \texttt{vnb-commands.lisp}
\item \texttt{vnb-complete.el}
\item \texttt{vnb-launch.el}
\item \texttt{vnb-panel-check.el}
\item \texttt{vnb.el}
\end{itemize}
The \texttt{reference/} directory is entirely \emph{generated} on prover load (\texttt{STRUCTURE-INDEX}, \texttt{TACTICS}, \texttt{OPERATORS}, \texttt{THEOREMS}, \texttt{BY-TOPIC}, \texttt{PROOF-DEBT}, \ldots) as \texttt{.md}, converted to the browser HTML by \texttt{build-reference-html} (reference-html.scm); it is never hand-edited.


%% =========================================================
